\documentclass[11pt]{article}
\usepackage{amsmath}
\usepackage[margin=1in]{geometry}

\title{Sustained Space Maneuver: Technical Notes}
\author{}
\date{}

\begin{document}
\maketitle

\section{Scope and convention}

These notes retain mathematical and system-level material relevant to sustained
space maneuver. Unless stated otherwise, orbit descriptions use an
Earth-centered inertial frame and an explicitly stated epoch.

\section{Delta-v}

Delta-v, denoted $\Delta v$, is a measure of the velocity change supplied by a
spacecraft's propulsion system. It is reported in units of speed, commonly
$\mathrm{m/s}$ or $\mathrm{km/s}$, but is principally a measure of maneuvering
capability and propellant demand.

For an impulsive maneuver, the velocity-vector change is
\begin{equation}
  \Delta \boldsymbol{v} = \boldsymbol{v}^{+} - \boldsymbol{v}^{-},
\end{equation}
where $-$ and $+$ denote the states immediately before and after the burn. The
associated scalar delta-v is $\lVert\Delta\boldsymbol{v}\rVert$.

For a finite, sustained thrust maneuver, an idealized propulsive delta-v
expenditure is
\begin{equation}
  \Delta v_{\mathrm{prop}} = \int_{t_0}^{t_f} \lVert\boldsymbol{a}_{T}(t)\rVert\,dt,
\end{equation}
where $\boldsymbol{a}_{T}$ is the thrust acceleration. This is distinct from the
net inertial velocity change, because gravity and other perturbing accelerations
also change the spacecraft velocity.

The ideal rocket equation relates available delta-v to the mass ratio:
\begin{equation}
  \Delta v = I_{sp} g_0 \ln\!\left(\frac{m_0}{m_f}\right),
\end{equation}
where $I_{sp}$ is specific impulse, $g_0$ is standard gravity, and $m_0$ and
$m_f$ are the initial and final vehicle masses. A mission delta-v budget is the
cumulative propulsive delta-v allocated across maneuvers, including any
station-keeping or sustained low-thrust operations.

\section{Classical orbital elements}

The six classical orbital elements (COEs) specify a Keplerian orbit at a given
epoch:
\begin{center}
\begin{tabular}{lll}
  $a$ & semi-major axis & orbit size and specific orbital energy \\
  $e$ & eccentricity & orbit shape \\
  $i$ & inclination & tilt of the orbital plane \\
  $\Omega$ & right ascension of the ascending node (RAAN) & orientation of the line of nodes \\
  $\omega$ & argument of periapsis & orientation of periapsis within the plane \\
  $\nu$ & true anomaly & instantaneous position measured from periapsis
\end{tabular}
\end{center}

Here, $\Omega$ is measured in the reference plane from the positive $x$ axis to
the ascending node; $\omega$ is measured in the orbital plane from the ascending
node to periapsis; and $\nu$ is measured in the orbital plane from periapsis to
the spacecraft. The usual element set has singularities for circular or
equatorial orbits: periapsis is undefined when $e=0$, and the line of nodes is
undefined when $i=0$ or $\pi$. Alternative nonsingular element sets should be
used in those regimes.

The right ascension of the ascending node (RAAN), $\Omega$, therefore specifies
the inertial direction of the orbit plane's northbound equatorial crossing. In a
J2000 ECI frame it is measured eastward, from the $+x$ direction to the
ascending node, over $[0,2\pi)$. RAAN identifies the line of nodes rather than
the spacecraft's instantaneous location on the orbit.

Perturbations cause osculating COEs to vary with time. Consequently, a reported
element set should state the reference frame, the epoch, and whether the elements
are osculating or mean elements.

\section{J2000 Earth-centered inertial frame}

J2000 ECI denotes an Earth-centered, non-Earth-rotating reference frame whose
orientation is tied to the J2000 epoch. Its origin is Earth's center of mass.
Under the conventional J2000/EME2000 description, the $+z$ axis is aligned with
the mean terrestrial north-pole direction at J2000, the $+x$ axis points toward
the mean equinox of J2000, and the $+y$ axis completes a right-handed triad. The
J2000.0 epoch is Julian Date $2451545.0$, corresponding to 2000 January 1,
12:00~TT.

An ECI frame does not rotate daily with Earth and is therefore convenient for
propagation and expressing inertial position $\boldsymbol r$ and velocity
$\boldsymbol v$. An Earth-fixed frame instead rotates with Earth and is more
appropriate for ground tracks and terrestrial locations. In precise work,
EME2000 and the Geocentric Celestial Reference Frame (GCRF) should not be
treated as exactly interchangeable, even though the distinction is often small
for introductory or approximate applications.

\section{Point of closest approach}

For two objects with inertial position vectors $\boldsymbol r_1(t)$ and
$\boldsymbol r_2(t)$, define the relative position and separation as
\begin{equation}
  \boldsymbol\rho(t)=\boldsymbol r_2(t)-\boldsymbol r_1(t),
  \qquad d(t)=\lVert\boldsymbol\rho(t)\rVert.
\end{equation}
Over a specified encounter window $[t_0,t_f]$, the time of closest approach is
\begin{equation}
  t_{\mathrm{TCA}}=\underset{t\in[t_0,t_f]}{\operatorname{arg\,min}}\;d(t),
\end{equation}
and the corresponding miss distance is $d_{\min}=d(t_{\mathrm{TCA}})$. The point
of closest approach (POCA) refers to the state or location associated with this
minimum-separation event. Because each object occupies a different inertial
position at TCA, POCA should be qualified as the point on a particular object's
trajectory or as a point on the relative trajectory.

At a smooth interior minimum, the relative position and relative velocity are
orthogonal:
\begin{equation}
  \boldsymbol\rho(t_{\mathrm{TCA}})\mathbin{\cdot}
  \dot{\boldsymbol\rho}(t_{\mathrm{TCA}})=0.
\end{equation}
For sustained maneuvering, POCA and TCA are control-dependent predictions rather
than fixed properties of the original trajectories. Continuous thrust can move,
remove, or create local minima in separation, so encounter analysis should state
the prediction window and distinguish the global minimum from local minima.

\section{TDA-informed static refueling-depot placement}

Let $c_{ij}$ denote the transfer cost, expressed initially as delta-v, from
customer orbit $i$ to candidate depot orbit $j$. For a selected depot set
$S$, order the customer-specific costs as
\begin{equation}
  c_{i(1)}(S)\leq c_{i(2)}(S)\leq\cdots.
\end{equation}
The primary service cost is $c_{i(1)}$, while $c_{i(q)}$ measures the cost of
accessing the $q$th depot and therefore represents $q$-fold service redundancy.
Coverage at an allowable cost $\epsilon_*$ requires
\begin{equation}
  \left|\{j\in S:c_{ij}\leq\epsilon_*\}\right|\geq q
\end{equation}
for each required customer orbit.

A topological filtration is obtained by constructing, for each $\epsilon$, a
bipartite graph whose vertices are customer orbits and selected depots, with an
edge $(i,j)$ whenever $c_{ij}\leq\epsilon$. The zeroth Betti number
$\beta_0(\epsilon)$ counts connected components and exposes fragmentation. For
the graph filtration, the first Betti number is
\begin{equation}
  \beta_1(\epsilon)=|E|-|V|+\beta_0(\epsilon),
\end{equation}
which counts independent cycles. In a bipartite service graph, cycles indicate
alternate customer--depot paths and can be used as one interpretable redundancy
signal. Explicit degree-based $q$-coverage and failure analysis remain necessary;
$\beta_1$ alone does not certify operational redundancy.

A baseline multi-criteria objective may combine weighted mean primary delta-v,
the $q$th-depot delta-v, worst cost after a single depot failure, coverage
deficit, and topological fragmentation. The terms should be normalized or
reported as a Pareto frontier rather than combined with arbitrary coefficients.
The transfer-cost model is a separate layer and should ultimately include
epoch, phasing, time of flight, finite thrust, propellant delivery, replenishment,
and other operational constraints.

\subsection{Dynamic service network}

A dynamic implementation replaces the static cost matrix with a directed,
time-dependent set of transfer opportunities. A transfer arc may be written
\begin{equation}
  a=(i,j,t_{\mathrm{dep}},t_{\mathrm{arr}},m_0,v),
\end{equation}
where $m_0$ is initial vehicle mass and $v$ identifies the vehicle or propulsion
model. Each arc carries a vector of costs and attributes, including delta-v,
propellant mass, time of flight, and probability of success. Such costs are not
generally symmetric and need not satisfy the triangle inequality.

The operational system can be represented as a directed time-expanded network.
For depot $j$, a basic inventory balance is
\begin{equation}
  I_{j,t+1}=I_{j,t}+D_{j,t}-S_{j,t}-L_{j,t},
\end{equation}
where $D$ is delivered propellant, $S$ is serviced demand, and $L$ represents
boil-off and other losses. Redundancy should distinguish geometric $q$-coverage,
multiple time-respecting service paths, and sufficient inventory and vehicle
capacity along those paths.

Because service complexes can both gain and lose edges over time, ordinary
nested persistence does not directly describe their temporal evolution. Snapshot
persistence provides a baseline; zigzag persistence can represent additions and
deletions; and persistent path homology is a candidate for retaining directed
transfer asymmetry. These topological summaries should be compared against
standard graph measures and tested under uncertainty before inclusion in an
optimization objective.

\end{document}
